\documentclass[12pt,a4paper]{article}

% ── Packages ──────────────────────────────────────────────────
\usepackage[utf8]{inputenc}
\usepackage[T1]{fontenc}
% Language settings
\usepackage{babel}

\usepackage{geometry}
\usepackage{amsmath}
\usepackage{amssymb}
\usepackage{hyperref}
\usepackage{xcolor}
\usepackage{listings}
\usepackage{booktabs}
\usepackage{longtable}
\usepackage{array}
\usepackage{graphicx}
\usepackage{fancyhdr}
\usepackage{titlesec}
\usepackage{tcolorbox}
\usepackage{enumitem}
\usepackage{parskip}

% ── Page geometry ─────────────────────────────────────────────
\geometry{
  headheight=25pt,
  a4paper,
  top=2.5cm, bottom=2.5cm,
  left=2.5cm, right=2.5cm
}

% ── Colors ────────────────────────────────────────────────────
\definecolor{bleu}{RGB}{31,56,100}
\definecolor{bleu2}{RGB}{46,116,181}
\definecolor{bleu3}{RGB}{212,239,239}
\definecolor{gris}{RGB}{245,245,245}
\definecolor{gris2}{RGB}{200,200,200}
\definecolor{vert}{RGB}{0,128,0}
\definecolor{rouge}{RGB}{180,30,30}
\definecolor{violet}{RGB}{74,20,140}
\definecolor{orange}{RGB}{197,90,17}

% ── Hyperref ──────────────────────────────────────────────────
\hypersetup{
  colorlinks=true,
  linkcolor=bleu2,
  urlcolor=bleu2,
  citecolor=bleu2,
  pdftitle={SICoxCURE -- User Guide v0.1.2},
  pdfauthor={Mohamed Yassir Errahmani}
}

% ── Headers / Footers ─────────────────────────────────────────
\pagestyle{fancy}
\fancyhf{}
\fancyhead[L]{\textcolor{bleu}{\textbf{SICoxCURE}}}
\fancyhead[R]{\textcolor{bleu}{v0.1.2 -- User Guide}}
\fancyfoot[C]{\textcolor{gris2}{\thepage}}
\renewcommand{\headrulewidth}{1pt}
\renewcommand{\headrule}{\hbox to\headwidth{\color{bleu}\leaders\hrule height \headrulewidth\hfill}}

% ── Section style ─────────────────────────────────────────────
\titleformat{\section}
  {\Large\bfseries\color{bleu}}
  {\thesection.}{0.5em}{}
  [\vspace{-0.3em}\textcolor{bleu2}{\rule{\linewidth}{1pt}}]

\titleformat{\subsection}
  {\large\bfseries\color{bleu2}}
  {\thesubsection}{0.5em}{}

\titleformat{\subsubsection}
  {\normalsize\bfseries\color{orange}}
  {\thesubsubsection}{0.5em}{}

% ── R code listings ───────────────────────────────────────────
\lstdefinestyle{Rcode}{
  language=R,
  basicstyle=\ttfamily\small\color{black},
  keywordstyle=\color{bleu2}\bfseries,
  stringstyle=\color{rouge},
  commentstyle=\color{vert}\itshape,
  numberstyle=\tiny\color{gris2},
  numbers=left,
  numbersep=8pt,
  backgroundcolor=\color{gris},
  frame=single,
  framerule=1pt,
  rulecolor=\color{bleu2},
  breaklines=true,
  breakatwhitespace=true,
  showstringspaces=false,
  tabsize=2,
  captionpos=b,
  xleftmargin=15pt,
  morekeywords={sic_fit, bootstrap_se, plot_link,
                plot_survival, predict, summary,
                print, library, matrix, rbinom,
                rexp, rnorm, pmin, ifelse, as.integer}
}

% Style pour les sorties R
\lstdefinestyle{Routput}{
  basicstyle=\ttfamily\footnotesize\color{black},
  backgroundcolor=\color{white},
  frame=single,
  framerule=0.5pt,
  rulecolor=\color{gris2},
  breaklines=true,
  xleftmargin=15pt,
  numbers=none,
}

% ── Boites colorées ───────────────────────────────────────────
\tcbuselibrary{skins,breakable}

\newtcolorbox{infobox}[1][]{
  colback=bleu3, colframe=bleu2,
  fonttitle=\bfseries\color{bleu},
  title=#1, breakable,
  left=6pt, right=6pt, top=4pt, bottom=4pt
}

\newtcolorbox{warningbox}[1][]{
  colback=yellow!10, colframe=orange,
  fonttitle=\bfseries\color{orange},
  title=#1, breakable,
  left=6pt, right=6pt, top=4pt, bottom=4pt
}

\newtcolorbox{formulebox}{
  colback=bleu3!50, colframe=bleu,
  center, breakable,
  left=10pt, right=10pt, top=8pt, bottom=8pt
}

% ══════════════════════════════════════════════════════════════
\begin{document}
% ══════════════════════════════════════════════════════════════

% ── TITLE PAGE ────────────────────────────────────────────────
\begin{titlepage}
  \centering
  \vspace*{2cm}

  {\color{bleu}\rule{\linewidth}{2pt}}
  \vspace{0.5cm}

  {\Huge\bfseries\color{bleu} SICoxCURE}\\[0.3cm]
  {\Large\color{bleu2} Single-Index/Cox Mixture Cure Model}\\[0.2cm]
  {\large\color{black} R Package --- Version 0.1.2}

  \vspace{0.3cm}
  {\color{bleu}\rule{\linewidth}{2pt}}

  \vspace{1.5cm}

  {\large\textbf{User Guide \& Reference Manual}}

  \vspace{2cm}

  \begin{center}
\begin{tabular}{rl}
    \textbf{Developer:}   & Mohamed Yassir Errahmani \\[0.2cm]
    \textbf{GitHub:}      & \href{https://github.com/MohamedYassirErrahmani/SICoxCURE}
                            {github.com/MohamedYassirErrahmani/SICoxCURE} \\[0.5cm]
    \textbf{Authors:}     & Mohamed Yassir Errahmani \& Maïlis Amico \\[0.2cm]
    \textbf{Email:}       & \href{mailto:mohamed-yassir.errahmani@umontpellier.fr}
                            {mohamed-yassir.errahmani@umontpellier.fr} \\
                          & \href{mailto:mailis.amico@kuleuven.be}
                            {mailis.amico@kuleuven.be} \\[0.5cm]
    \textbf{Reference:}   & Amico, Van Keilegom \& Legrand (2019) \\
                          & \textit{Biometrics}, \textbf{75}(2), 452--462 \\
                          & \href{https://doi.org/10.1111/biom.12999}
                            {doi:10.1111/biom.12999} \\[0.5cm]
    \textbf{License:}     & GPL ($\geq$ 3)
  \end{tabular}
  \end{center}

  \vfill

  {\color{bleu}\rule{\linewidth}{1pt}}
  {\small\color{gris2} SICoxCURE v0.1.2 ---
   \href{https://github.com/MohamedYassirErrahmani/SICoxCURE}
   {github.com/MohamedYassirErrahmani/SICoxCURE}}
\end{titlepage}

% ── TABLE OF CONTENTS ─────────────────────────────────────────
\tableofcontents
\newpage

% ══════════════════════════════════════════════════════════════
\section{Overview}
% ══════════════════════════════════════════════════════════════

The \textbf{SICoxCURE} package implements the
\textbf{Single-Index/Cox (SIC) mixture cure model} proposed by
Amico, Van Keilegom and Legrand (2019). It was developed and
maintained by \textbf{Mohamed Yassir Errahmani}
(University of Montpellier).

\subsection{The Model}

In survival analysis, a cure model is appropriate when a fraction
of subjects will never experience the event of interest (the
``cured'' subjects). The SIC mixture cure model writes the
population survival function as:

\begin{formulebox}
  \[
    S(t \mid \mathbf{x}, \mathbf{z}) =
    \underbrace{1 - p(\mathbf{x})}_{\text{cured fraction}}
    +
    \underbrace{p(\mathbf{x}) \cdot S_u(t \mid \mathbf{z})}_{
      \text{uncured fraction}}
  \]
\end{formulebox}

\medskip

where:
\begin{itemize}[leftmargin=2cm]
  \item $p(\mathbf{x}) = g(\boldsymbol{\gamma}^\top \mathbf{x})$
    is the \textbf{incidence} (probability of being uncured),
    modeled via a single-index with unknown link function $g$
    estimated nonparametrically using Nadaraya-Watson kernel
    smoothing.
  \item $S_u(t \mid \mathbf{z})$ is the \textbf{latency}
    (conditional survival of uncured subjects), modeled via a
    Cox proportional hazards model.
  \item The vectors $\mathbf{X}$ and $\mathbf{Z}$ can contain
    partially or completely different covariates.
\end{itemize}

\subsection{Estimation}

Estimation is performed by maximum likelihood using an
\textbf{Expectation-Maximization (EM) algorithm}:
\begin{itemize}
  \item \textbf{E-step:} computes posterior weights
    $W_i = P(B_i = 1 \mid \text{data},\, \theta^{(m-1)})$
  \item \textbf{M-step (incidence):} maximizes $\tilde{L}_1$
    to update $\hat{\boldsymbol{\gamma}}$ and
    $\hat{g}(\cdot)$ (Nadaraya-Watson estimator)
  \item \textbf{M-step (latency):} maximizes $\tilde{L}_2$
    (partial likelihood) to update $\hat{\boldsymbol{\beta}}$
    and $\hat{\Lambda}_0$
\end{itemize}

The bandwidth $h$ for the kernel estimator is selected at each
iteration by \textbf{likelihood cross-validation}
over a user-specified grid.

\subsection{When to Use a Cure Model}

\begin{infobox}[Conditions for a cure model]
\begin{enumerate}
  \item Event rate $< 50\%$ (many potential long-term survivors)
  \item Visible plateau in the Kaplan-Meier estimator ($> 10\%$)
  \item Sufficient follow-up time beyond the last observed event
\end{enumerate}
\end{infobox}

% ══════════════════════════════════════════════════════════════
\section{Installation}
% ══════════════════════════════════════════════════════════════

\subsection{From CRAN (Coming soon...)}

\begin{lstlisting}[style=Rcode, caption={Installation from CRAN}]
install.packages("SICoxCURE")
\end{lstlisting}

\subsection{From GitHub (development version)}

\begin{lstlisting}[style=Rcode, caption={Installation from GitHub}]
# Install devtools if needed
install.packages("devtools")

# Install SICoxCURE from GitHub
devtools::install_github("MohamedYassirErrahmani/SICoxCURE")
\end{lstlisting}

\subsection{Load the Package}

\begin{lstlisting}[style=Rcode, caption={Loading SICoxCURE}]
library(SICoxCURE)

# Check version
packageVersion("SICoxCURE")
\end{lstlisting}

\begin{lstlisting}[style=Routput, caption={Expected output}]
[1] '0.1.2'
\end{lstlisting}

\begin{lstlisting}[style=Rcode, caption={List available functions}]
ls("package:SICoxCURE")
\end{lstlisting}

\begin{lstlisting}[style=Routput]
[1] "bootstrap_se"  "plot_link"     "plot_survival" "sic_fit"
\end{lstlisting}

\begin{lstlisting}[style=Rcode, caption={Check S3 methods}]
methods(class = "sic_fit")
\end{lstlisting}

\begin{lstlisting}[style=Routput]
[1] predict   print   summary
\end{lstlisting}

% ══════════════════════════════════════════════════════════════
\section{Data Preparation}
% ══════════════════════════════════════════════════════════════

\subsection{Required Data Structure}

The \texttt{sic\_fit()} function requires the following inputs:

\begin{center}
\begin{tabular}{lll}
  \toprule
  \textbf{Argument} & \textbf{Type} & \textbf{Description} \\
  \midrule
  \texttt{Y}     & Numeric vector $(n)$ &
    Follow-up times $Y_i = \min(T_i, C_i)$ \\
  \texttt{delta} & Integer vector $(n)$ &
    Event indicators ($1$=event, $0$=censored) \\
  \texttt{X}     & Matrix $(n \times d)$ &
    Incidence covariates for $\boldsymbol{\gamma}^\top\mathbf{X}$ \\
  \texttt{Z}     & Matrix $(n \times q)$ &
    Latency covariates for Cox model \\
  \bottomrule
\end{tabular}
\end{center}

\subsection{How to Organize Your Own Data}

Before fitting the model, the user must prepare three
objects:

\begin{itemize}
  \item \textbf{Y}: a numeric vector of follow-up times
    $Y_i = \min(T_i, C_i)$, one value per subject.
  \item \textbf{delta}: an integer vector of event
    indicators: $1$ if the event was observed,
    $0$ if censored.
  \item \textbf{X}: a numeric matrix with $n$ rows and
    $d$ columns, where each column is one incidence
    covariate (used to model the uncure probability).
  \item \textbf{Z}: a numeric matrix with $n$ rows and
    $q$ columns, where each column is one latency
    covariate (used in the Cox model for uncured subjects).
\end{itemize}

\begin{infobox}[Important]
\begin{itemize}
  \item \textbf{X} and \textbf{Z} can contain the same
    covariates, partially overlapping ones, or completely
    different ones.
  \item Covariates in \textbf{X} are automatically
    standardized (centered and scaled) by
    \texttt{sic\_fit()} when \texttt{standardize\_X = TRUE}
    (default). This follows the recommendation of
    Amico, Van Keilegom and Legrand (2019)
    \href{https://doi.org/10.1111/biom.12999}
    {doi:10.1111/biom.12999}.
  \item Missing values must be removed before calling
    \texttt{sic\_fit()}.
\end{itemize}
\end{infobox}

\textbf{Example with a real dataset:}

\begin{lstlisting}[style=Rcode,
  caption={Preparing your own data}]
# Load your data
dat <- read.csv("your_data.csv")

# Remove missing values
dat <- na.omit(dat)

# Prepare vectors and matrices
Y     <- dat$follow_up_time    # follow-up times
delta <- dat$event_indicator   # 1=event, 0=censored

# Incidence covariates (X)
X <- as.matrix(dat[, c("age", "ER_status",
                        "tumor_size")])

# Latency covariate (Z)
Z <- as.matrix(dat[, "ER_status", drop = FALSE])

# Check dimensions
cat(sprintf("n=%d | Events=%.1f%% | Censored=%.1f%%\n",
            nrow(dat),
            100 * mean(delta),
            100 * mean(delta == 0)))
\end{lstlisting}

\subsection{Using a data.frame}

\begin{lstlisting}[style=Rcode,
  caption={data.frame interface}]
dat <- data.frame(Y=Y, delta=delta,
                  X1=X[,1], X2=X[,2],
                  X3=X[,3], X4=X[,4],
                  Z=Z[,1])

fit <- sic_fit(
  data      = dat,
  col_Y     = "Y",
  col_delta = "delta",
  cols_X    = c("X1","X2","X3","X4"),
  cols_Z    = "Z"
)
\end{lstlisting}

% ══════════════════════════════════════════════════════════════
\section{Fitting the Model: \texttt{sic\_fit()}}
% ══════════════════════════════════════════════════════════════

\subsection{Main Parameters}

\begin{longtable}{lll}
  \toprule
  \textbf{Parameter} & \textbf{Default} & \textbf{Description} \\
  \midrule
  \texttt{Y}            & ---         & Follow-up times \\
  \texttt{delta}        & ---         & Event indicators \\
  \texttt{X}            & ---         & Incidence covariates \\
  \texttt{Z}            & ---         & Latency covariates \\
  \texttt{kernel}       & \texttt{"epanechnikov"} &
    Kernel function \\
  \texttt{h\_range}     & \texttt{c(0.4, 1.0)} &
    Bandwidth search interval \\
  \texttt{h\_step}      & \texttt{0.001} &
    Bandwidth grid step (601 values) \\
  \texttt{constraint}   & \texttt{"gamma1"} &
    Identifiability: fix $\gamma_1=1$ \\
  \texttt{tol}          & \texttt{1e-5} &
    EM convergence tolerance \\
  \texttt{max\_iter}    & \texttt{500L} &
    Maximum EM iterations \\
  \texttt{standardize\_X} & \texttt{TRUE} &
    Standardize X before fitting \\
  \texttt{bootstrap}    & \texttt{FALSE} &
    Compute SE automatically \\
  \texttt{B}            & \texttt{250L} &
    Bootstrap samples \\
  \texttt{boot\_seed}   & \texttt{NULL} &
    Seed for bootstrap \\
  \texttt{plot}         & \texttt{FALSE} &
    Auto-plot after fitting \\
  \texttt{plot\_Z}      & \texttt{NULL} &
    Z values for survival plot \\
  \texttt{plot\_labels} & \texttt{c("Z=0","Z=1")} &
    Labels for survival plot \\
  \texttt{verbose}      & \texttt{1L} &
    0=silent, 1=progress, 2=detail \\
  \bottomrule
\end{longtable}

\subsection{Workflow A: All-in-One (Recommended)}

\begin{lstlisting}[style=Rcode,
  caption={Fitting with integrated bootstrap and plots}]
fit <- sic_fit(
  Y           = Y,
  delta       = delta,
  X           = X,
  Z           = Z,
  h_step      = 0.05,        # coarser grid for speed
  bootstrap   = TRUE,        # compute SE automatically
  B           = 250L,        # 250 bootstrap samples (article)
  boot_seed   = 42L,         # for reproducibility
  plot        = TRUE,        # auto-plot after fitting
  plot_Z      = matrix(c(0, 1), ncol = 1),
  plot_labels = c("Z = 0", "Z = 1"),
  verbose     = 1L           # show progress bar
)
\end{lstlisting}

\textbf{Console output (observed during fitting):}

\begin{lstlisting}[style=Routput]
=== Fitting SIC Mixture Cure Model ===
Step 1/3: Data validation and preparation...
X standardized (centered and scaled): 4 column(s). [Article]
  X1: mean=-0.010, sd=1.010 -> after: mean=0, sd=1
  ...
=== Data summary ===
  n=200 (no missing data)
  Events=89 (44.5%) | Censored=111
  tau=3.4521 | Plateau=48 (24.0%)
  X [X1,X2,X3,X4] | Z [Z]
====================
Step 2/3: Initialisation...
Step 3/3: EM algorithm...
--- Starting EM algorithm ---
Bandwidth grid: [0.40, 1.00] step=0.0500 (13 values)
  |===================================| 100%
[32m Checkmark Model converged successfully at iteration 8
        (criterion = 3.21e-06 < tol = 1e-05)[0m
  Final complete log-likelihood : -412.3847
Computing bootstrap standard errors...
Bootstrap standard errors: B=250, n=200, parallel=TRUE
\end{lstlisting}

\begin{infobox}[Convergence message]
When the model converges, a green message is displayed:
\textcolor{vert}{\texttt{\checkmark Model converged successfully at iteration N}}
along with the convergence criterion value.
\end{infobox}

\subsection{Workflow B: Step by Step}

\begin{lstlisting}[style=Rcode,
  caption={Step-by-step workflow}]
# Step 1: Fit the model
fit <- sic_fit(Y, delta, X, Z,
               h_step  = 0.05,
               verbose = 1L)

# Step 2: Bootstrap standard errors
se <- bootstrap_se(fit,
                   B        = 250L,
                   parallel = TRUE,
                   seed     = 42L,
                   verbose  = TRUE)

# Step 3: Summary with p-values
summary(fit, se = se)

# Step 4: Graphical outputs
plot_link(fit)
plot_survival(fit,
              newZ         = matrix(c(0, 1), ncol = 1),
              group_labels = c("Z = 0", "Z = 1"))

# Step 5: Predictions
pred <- predict(fit,
                newX = X[1:5, ],
                newZ = Z[1:5, , drop=FALSE],
                type = "all")
\end{lstlisting}

% ══════════════════════════════════════════════════════════════
\section{Output Object: \texttt{sic\_fit}}
% ══════════════════════════════════════════════════════════════

The \texttt{sic\_fit()} function returns an S3 object of class
\texttt{"sic\_fit"} containing:

\begin{center}
\begin{tabular}{lll}
  \toprule
  \textbf{Element} & \textbf{Type} & \textbf{Description} \\
  \midrule
  \texttt{gamma}     & Vector $(d)$ &
    Estimated incidence coefficients $\hat{\boldsymbol{\gamma}}$ \\
  \texttt{beta}      & Vector $(q)$ &
    Estimated latency coefficients $\hat{\boldsymbol{\beta}}$ \\
  \texttt{p\_hat}    & Vector $(n)$ &
    Estimated uncure probabilities $\hat{p}(\mathbf{x}_i)$ \\
  \texttt{Su\_hat}   & Vector $(n)$ &
    Estimated latency $\hat{S}_u(Y_i \mid \mathbf{z}_i)$ \\
  \texttt{breslow}   & List &
    Breslow estimator: times, $\Lambda_0$, $S_0$ \\
  \texttt{h}         & Scalar &
    Selected bandwidth $h_{\mathrm{CV}}$ \\
  \texttt{W}         & Vector $(n)$ &
    Final EM weights $\hat{W}_i$ \\
  \texttt{loglik}    & Vector &
    Complete log-likelihood per iteration \\
  \texttt{converged} & Logical &
    Did the EM algorithm converge? \\
  \texttt{n\_iter}   & Integer &
    Number of EM iterations \\
  \texttt{se}        & List &
    Bootstrap SE (if \texttt{bootstrap=TRUE}) \\
  \bottomrule
\end{tabular}
\end{center}

\subsection{print() method}

\begin{lstlisting}[style=Rcode, caption={Compact display}]
print(fit)
\end{lstlisting}

\begin{lstlisting}[style=Routput]
=== SIC Mixture Cure Model ===

Call: sic_fit(Y = Y, delta = delta, X = X, Z = Z, h_step = 0.05,
    verbose = 1L)

n = 200 | events = 89 | kernel = epanechnikov | h = 0.5500
Convergence: YES in 8 iterations (crit = 3.21e-06)

Incidence coefficients (gamma):
    X1     X2     X3     X4
1.0000 -0.4832  0.2876 -0.1954

Latency coefficients (beta):
    Z
1.3847

Use summary() for standard errors and p-values.
\end{lstlisting}

\subsection{summary() method}

\begin{lstlisting}[style=Rcode, caption={Summary with p-values}]
summary(fit)   # uses fit$se if bootstrap=TRUE
# or
summary(fit, se = se)   # with external bootstrap SE
\end{lstlisting}

\begin{lstlisting}[style=Routput]
=== Summary: SIC Mixture Cure Model ===

Call: sic_fit(Y = Y, delta = delta, X = X, Z = Z, ...)

Data:
  n = 200 | Events: 89 (44.5%) | Cure threshold tau = 3.4521
  Plateau: 48 observations (24.0%)

Model:
  Kernel: epanechnikov | Bandwidth h = 0.5500
  Constraint: 'gamma1' | Standardized X: yes

Convergence: YES in 8 iterations (crit = 3.21e-06, tol = 1e-05)

Incidence (gamma) -- single-index coefficients:
   Estimate Std.Error z.value Pr(>|z|)
X1   1.0000        --      --       --
X2  -0.4832    0.1943  -2.487   0.0129  *
X3   0.2876    0.2341   1.229   0.2192
X4  -0.1954    0.1876  -1.042   0.2974

Latency (beta) -- Cox model coefficients:
   Estimate Std.Error z.value Pr(>|z|)
Z    1.3847    0.2134   6.489   <0.001 ***

Signif. codes: '<0.001' '***' 0.001 '**' 0.01 '*' 0.05 '.' 0.1
\end{lstlisting}

\begin{infobox}[p-values]
P-values below 0.001 are displayed as \texttt{<0.001} instead
of \texttt{0}. Standard errors require bootstrap (either via
\texttt{bootstrap=TRUE} in \texttt{sic\_fit()} or via a
separate call to \texttt{bootstrap\_se()}).
\end{infobox}

% ══════════════════════════════════════════════════════════════
\section{Bootstrap Standard Errors: \texttt{bootstrap\_se()}}
% ══════════════════════════════════════════════════════════════

\begin{lstlisting}[style=Rcode,
  caption={Computing bootstrap SE (as in the article, B=250)}]
se <- bootstrap_se(
  fit      = fit,
  B        = 250L,      # as in Amico et al. (2019)
  parallel = TRUE,      # use multiple CPU cores
  ncores   = NULL,      # auto: all cores minus 1
  seed     = 42L,       # reproducibility
  verbose  = TRUE
)
\end{lstlisting}

\begin{lstlisting}[style=Routput]
Bootstrap standard errors: B=250, n=200, parallel=TRUE
Using 7 cores
This may take several minutes...
  |===================================| 100%

Done. 250/250 iterations succeeded.
\end{lstlisting}

\begin{lstlisting}[style=Rcode, caption={Display bootstrap SE}]
print(se)
\end{lstlisting}

\begin{lstlisting}[style=Routput]
=== Bootstrap Standard Errors (SICoxCURE) ===

Bootstrap samples : 250 / 250 succeeded

Incidence (gamma):
    X1     X2     X3     X4
0.0000 0.1943 0.2341 0.1876

Latency (beta):
    Z
0.2134

Use summary(fit, se = <this object>) for p-values.
\end{lstlisting}

% ══════════════════════════════════════════════════════════════
\section{Graphical Functions}
% ══════════════════════════════════════════════════════════════

\subsection{\texttt{plot\_link()}: Estimated Link Function}

\begin{lstlisting}[style=Rcode,
  caption={Plotting the estimated link function g}]
plot_link(fit,
          main = "Estimated link function g",
          col  = "blue",
          lwd  = 2)
\end{lstlisting}

This function plots $\hat{g}(u)$ against the single index
$\hat{\boldsymbol{\gamma}}^\top \mathbf{x}_i$ for the training
observations, sorted by index value. A horizontal dashed line
indicates the mean uncure probability.

\subsection{\texttt{plot\_survival()}: Latency Survival Curves}

\begin{lstlisting}[style=Rcode,
  caption={Plotting latency survival S\_u(t|z)}]
plot_survival(
  fit,
  newZ         = matrix(c(0, 1), ncol = 1),
  group_labels = c("Z = 0", "Z = 1"),
  cols         = c("red", "blue"),
  main         = "Estimated latency survival"
)
\end{lstlisting}

This function plots $\hat{S}_u(t \mid \mathbf{z})$ for specified
values of \texttt{Z}. It uses the Breslow estimator stored in
\texttt{fit\$breslow} and the estimated $\hat{\boldsymbol{\beta}}$.

\subsection{Integrated in \texttt{sic\_fit()}}

Both graphical functions can be called automatically by setting
\texttt{plot = TRUE} in \texttt{sic\_fit()}:

\begin{lstlisting}[style=Rcode,
  caption={Automatic plots within sic\_fit()}]
fit <- sic_fit(Y, delta, X, Z,
               plot        = TRUE,
               plot_Z      = matrix(c(0, 1), ncol = 1),
               plot_labels = c("Z = 0", "Z = 1"))
\end{lstlisting}

% ══════════════════════════════════════════════════════════════
\section{Predictions: \texttt{predict()}}
% ══════════════════════════════════════════════════════════════

\begin{lstlisting}[style=Rcode,
  caption={Predicting for new observations}]
# Generate new observations
set.seed(99)
X_new <- matrix(rnorm(10 * 4), 10, 4)
Z_new <- matrix(rbinom(10, 1L, 0.6), 10, 1)

# Incidence only
p_hat <- predict(fit,
                 newX = X_new,
                 newZ = Z_new,
                 type = "incidence")

# All predictions
pred <- predict(fit,
                newX   = X_new,
                newZ   = Z_new,
                t_eval = c(0.5, 1, 2, 5),
                type   = "all")
\end{lstlisting}

\subsection{Components of \texttt{pred}}

\begin{center}
\begin{tabular}{lll}
  \toprule
  \textbf{Component} & \textbf{Type} & \textbf{Description} \\
  \midrule
  \texttt{pred\$incidence} & Vector $(m)$ &
    $\hat{p}(\mathbf{x}_i)$ for each patient \\
  \texttt{pred\$latency}   & Matrix $(m \times T)$ &
    $\hat{S}_u(t \mid \mathbf{z}_i)$ \\
  \texttt{pred\$survival}  & Matrix $(m \times T)$ &
    $\hat{S}(t \mid \mathbf{x}_i, \mathbf{z}_i)$ \\
  \texttt{pred\$t\_eval}   & Vector $(T)$ &
    Time points \\
  \texttt{pred\$index}     & Vector $(m)$ &
    $\hat{\boldsymbol{\gamma}}^\top \mathbf{x}_i$ \\
  \bottomrule
\end{tabular}
\end{center}

\begin{lstlisting}[style=Rcode,
  caption={Displaying predictions as a table}]
tab <- data.frame(
  Patient   = 1:10,
  Index     = round(pred$index, 4),
  p_x       = round(pred$incidence, 4),
  S_t1      = round(pred$survival[, 2], 4),   # t=1
  S_t5      = round(pred$survival[, 4], 4)    # t=5
)
colnames(tab) <- c("Patient","Index","p(x)",
                   "S(t=1|x,z)","S(t=5|x,z)")
print(tab)
\end{lstlisting}

\begin{lstlisting}[style=Routput]
   Patient   Index   p(x) S(t=1|x,z) S(t=5|x,z)
1        1  1.2138 0.9960      0.004      0.004
2        2  2.0503 1.0000      0.021      0.021
3        3 -0.9207 0.8973      0.103      0.103
4        4 -5.4874 0.0000      1.000      1.000
5        5 -8.7415 0.8374      0.163      0.163
6        6 -2.1682 0.4968      0.503      0.503
7        7 -2.0149 0.5450      0.455      0.455
8        8 -4.8327 0.0000      1.000      1.000
9        9 -4.0791 0.1621      0.838      0.838
10      10  0.5943 0.9911      0.030      0.030
\end{lstlisting}

\begin{infobox}[Interpreting predictions]
\begin{itemize}
  \item \textbf{Patient 4 / 8:} $p(x) \approx 0$ $\Rightarrow$
    cured $\Rightarrow$ $S(t|x,z) = 1.000$ \\
  \item \textbf{Patient 1 / 2:} $p(x) \approx 1$ $\Rightarrow$
    uncured $\Rightarrow$ poor prognosis \\
  \item \textbf{Patient 6:} $p(x) \approx 0.50$ $\Rightarrow$
    uncertain (50/50)
\end{itemize}
The formula $S = 1 - p(x) + p(x) \cdot S_u$ guarantees that
cured patients ($p \to 0$) always have $S \to 1$.
\end{infobox}

% ══════════════════════════════════════════════════════════════
\section{Real Data Example: VERIDEX}
% ══════════════════════════════════════════════════════════════

This example reproduces the breast cancer application of Amico
et al.\ (2019), Section 4. The dataset consists of 286 patients
with lymph-node-negative breast cancer (Wang et al., 2005).
The event of interest is distant metastasis.

\begin{lstlisting}[style=Rcode,
  caption={VERIDEX application (as in the article)}]
library(SICoxCURE)

# Load data
dat <- read.table("VERIDEX-survival.txt",
                  header = TRUE, sep = "\t")
colnames(dat) <- c("PID","GEO","lymph","time",
                   "status","ER","brain")
dat$ER_bin <- as.integer(dat$ER == "ER+")

cat(sprintf("n=%d | Events=%d (%.1f%%)\n",
            nrow(dat), sum(dat$status),
            100*mean(dat$status)))

# Train/test split (2/3 -- 1/3 as in the article)
set.seed(2019)
idx <- sample(nrow(dat), floor(2/3 * nrow(dat)))
train <- dat[idx, ]
X_tr  <- as.matrix(train[, "ER_bin", drop=FALSE])
Z_tr  <- as.matrix(train[, "ER_bin", drop=FALSE])

# Fit SICoxCURE (B=250 as in the article)
fit <- sic_fit(
  Y           = train$time,
  delta       = train$status,
  X           = X_tr,
  Z           = Z_tr,
  h_step      = 0.001,        # article default
  bootstrap   = TRUE,
  B           = 250L,
  boot_seed   = 2019L,
  plot        = TRUE,
  plot_Z      = matrix(c(0, 1), ncol = 1),
  plot_labels = c("ER- (Z=0)", "ER+ (Z=1)"),
  verbose     = 1L
)

# Summary (reproduces Table 3 of the article)
summary(fit)
\end{lstlisting}

\begin{lstlisting}[style=Routput]
n=286 | Events=107 (37.4%)

=== Summary: SIC Mixture Cure Model ===

Incidence (gamma):
          Estimate Std.Error z.value Pr(>|z|)
ER_bin      1.0000        --      --       --

Latency (beta):
          Estimate Std.Error z.value Pr(>|z|)
ER_bin     -0.9279    0.2753  -3.370  <0.001 ***

Signif. codes: '<0.001' '***' 0.001 '**' 0.01 '*' 0.05 '.' 0.1
\end{lstlisting}

% ══════════════════════════════════════════════════════════════
\section{Tips and Recommendations}
% ══════════════════════════════════════════════════════════════

\subsection{Bandwidth Selection}

\begin{warningbox}[Bandwidth at the boundary]
If $h_{\mathrm{CV}} = h_{\min}$ or $h_{\mathrm{CV}} = h_{\max}$,
the true optimum may lie outside the search interval. Extend
\texttt{h\_range} accordingly.
\end{warningbox}

\begin{lstlisting}[style=Rcode,
  caption={Bandwidth tuning}]
# Quick exploration (13 values)
fit_fast <- sic_fit(Y, delta, X, Z,
                    h_step = 0.05)

# Final publication result (601 values, as in article)
fit_final <- sic_fit(Y, delta, X, Z,
                     h_step = 0.001)
\end{lstlisting}

\subsection{Bootstrap}

\begin{lstlisting}[style=Rcode,
  caption={Bootstrap recommendations}]
# Quick exploration
se_quick <- bootstrap_se(fit, B = 30L,
                         parallel = FALSE)

# Publication quality (as in article)
se_final <- bootstrap_se(fit, B = 250L,
                         parallel = TRUE,
                         seed = 42L)
\end{lstlisting}

\subsection{Parallel Computing}

\begin{lstlisting}[style=Rcode,
  caption={Check available CPU cores}]
parallel::detectCores()  # total cores on your PC
\end{lstlisting}

\begin{infobox}[Parallel computing]
\texttt{parallel = TRUE} uses all available CPU cores minus one,
keeping one core free for the operating system. On a machine
with 8 cores, this reduces bootstrap time from $\sim$8 hours
($B=250$ sequential) to $\sim$1 hour.
\end{infobox}

% ══════════════════════════════════════════════════════════════
\section{Function Reference (v0.1.2)}
% ══════════════════════════════════════════════════════════════

\subsection{Public Functions}

\begin{longtable}{lll}
  \caption{Public functions in SICoxCURE v0.1.2}\\
  \toprule
  \textbf{Function} & \textbf{File} & \textbf{Description} \\
  \midrule
  \texttt{sic\_fit()} & \texttt{sic\_fit.R} &
    Main function. Fits the SIC model via EM. \\
  \texttt{bootstrap\_se()} & \texttt{bootstrap.R} &
    Bootstrap SE for $\hat{\boldsymbol{\gamma}}$
    and $\hat{\boldsymbol{\beta}}$. \\
  \texttt{plot\_link()} & \texttt{plots.R} &
    Plot the estimated link function $\hat{g}(u)$. \\
  \texttt{plot\_survival()} & \texttt{plots.R} &
    Plot latency survival $\hat{S}_u(t \mid \mathbf{z})$. \\
  \texttt{predict.sic\_fit()} & \texttt{predict.R} &
    Predict $\hat{p}(\mathbf{x})$, $\hat{S}_u$,
    $\hat{S}$ for new data. \\
  \texttt{print.sic\_fit()} & \texttt{sic\_fit.R} &
    Compact display of a \texttt{sic\_fit} object. \\
  \texttt{summary.sic\_fit()} & \texttt{sic\_fit.R} &
    Coefficient table with SE and p-values. \\
  \texttt{print.sic\_bootstrap()} & \texttt{bootstrap.R} &
    Display a \texttt{sic\_bootstrap} object. \\
  \bottomrule
\end{longtable}

\subsection{Internal Functions (accessible via \texttt{:::})}

\begin{longtable}{ll}
  \caption{Key internal functions in SICoxCURE v0.1.2}\\
  \toprule
  \textbf{File / Function} & \textbf{Role} \\
  \midrule
  \multicolumn{2}{l}{\textbf{em\_algorithm.R}} \\
  \texttt{run\_em()} &
    Main EM loop: E $\to$ M1 $\to$ M2 $\to$ M3 $\to$
    convergence check \\[3pt]
  \multicolumn{2}{l}{\textbf{estep.R}} \\
  \texttt{estep()} &
    E-step: computes $W_i = P(B_i=1 \mid \text{data}, \theta)$ \\[3pt]
  \multicolumn{2}{l}{\textbf{mstep\_incidence.R}} \\
  \texttt{mstep\_incidence()} &
    M-step for $\boldsymbol{\gamma}$:
    L-BFGS-B with Nelder-Mead fallback \\
  \texttt{neg\_loglik\_incidence()} &
    Objective function $-\log(\tilde{L}_1)$ for \texttt{optim()} \\[3pt]
  \multicolumn{2}{l}{\textbf{mstep\_latency.R}} \\
  \texttt{mstep\_latency()} &
    M-step for $\boldsymbol{\beta}$: partial likelihood \\
  \texttt{neg\_partial\_loglik()} &
    Objective function $-\log(\bar{L}_2)$ for \texttt{optim()} \\[3pt]
  \multicolumn{2}{l}{\textbf{nadaraya\_watson.R}} \\
  \texttt{nw\_loo()} &
    Leave-one-out Nadaraya-Watson estimator of $g$ \\
  \texttt{nw\_predict()} &
    Nadaraya-Watson prediction for new observations \\
  \texttt{select\_bandwidth()} &
    Cross-validation bandwidth selection \\[3pt]
  \multicolumn{2}{l}{\textbf{utils.R (17 functions)}} \\
  \texttt{kernel\_epanechnikov()} &
    Epanechnikov kernel
    $K(u) = 3(1-u^2/5)/(4\sqrt{5}) \cdot \mathbf{1}(u^2 \leq 5)$ \\
  \texttt{breslow\_estimator()} &
    Weighted Breslow estimator for $\hat{\Lambda}_0$ \\
  \texttt{eval\_Su()} &
    Evaluate $S_u(t \mid \mathbf{z})$ on a time grid \\
  \texttt{compute\_index()} &
    Single index $u_i = \boldsymbol{\gamma}^\top \mathbf{x}_i$ \\
  $\ldots$ & $\ldots$ \\
  \bottomrule
\end{longtable}

% ══════════════════════════════════════════════════════════════
\section{Reference}
% ══════════════════════════════════════════════════════════════

\textbf{Amico M, Van Keilegom I, Legrand C (2019).}
The single-index/Cox mixture cure model.
\textit{Biometrics}, \textbf{75}(2), 452--462.
\href{https://doi.org/10.1111/biom.12999}{doi:10.1111/biom.12999}

\bigskip

\textbf{Wang Y, Klijn JGM, Zhang Y, \textit{et al.} (2005).}
Gene-expression profiles to predict distant metastasis of
lymph-node-negative primary breast cancer.
\textit{The Lancet}, \textbf{365}, 671--679.

\bigskip

\textbf{Sy JP, Taylor JMG (2000).}
Estimation in a Cox proportional hazards cure model.
\textit{Biometrics}, \textbf{56}, 227--236.

\vfill

\begin{center}
  \textcolor{gris2}{\rule{0.8\linewidth}{0.5pt}}\\[0.3cm]
  {\small\color{gris2}
    SICoxCURE v0.1.2 ---
    Developer: Mohamed Yassir Errahmani ---
    \href{https://github.com/MohamedYassirErrahmani/SICoxCURE}
    {github.com/MohamedYassirErrahmani/SICoxCURE}
  }
\end{center}

\end{document}
